\section{Research Objectives}
\label{sec:research-objectives}

The primary objective of this research is to develop a precise localisation system for autonomous vineyard robots that enables accurate spatial alignment between the robot, the 3D models, and the environment.

The research pursues the following sub-objectives:

\begin{enumerate}

    \item Develop a LiDAR-inertial SLAM system that achieves sub-decimetre absolute trajectory error in GNSS-degraded vineyard environments, addressing the structural ambiguity of row-crop corridors.

    \item Design a precise-alignment pipeline that refines the SLAM pose to sub-centimetre accuracy by registering live sensor data against a 3D Gaussian Splatting reference model using point cloud registration.

    \item Evaluate two sensing modalities --- stereo vision and LiDAR --- for the precise-alignment stage, comparing their registration accuracy and failure modes under vineyard field conditions.

    \item Integrate both localisation stages into a hierarchical architecture that provides pose estimates suitable for downstream robotic manipulation, such that the corrected pose feeds directly into the pruning arm's path planner.

\end{enumerate}

\subsection{Contributions}
\label{subsec:contributions}

This thesis makes the following contributions:

\begin{enumerate}
    % TODO: Fill in contributions once results and discussion chapters are complete.
    % Each item should state an accomplished outcome with quantitative evidence where possible.
    % Suggested structure:
    %   1. A multi-sensor SLAM system for [context], achieving [metric].
    %   2. A precise-alignment pipeline that [does X], demonstrating [Y].
    %   3. An experimental evaluation showing [Z].
    \item \textbf{TODO:} Contribution 1 --- SLAM system contribution and headline result.
    \item \textbf{TODO:} Contribution 2 --- Precise-alignment pipeline contribution and headline result.
    \item \textbf{TODO:} Contribution 3 --- Experimental evaluation or additional contribution.
\end{enumerate}